\subsection{THE MEAN VALUE THEOREM FOR VECTOR FUNCTIONS}

THEOREM 2. Let \(\mathbf{f}\) be a vector function defined and continuous on a closed bounded interval \(\mathrm{I} = [a, b]\) of \(\mathbf{R}\) (where \(a < b\) ) and taking values in a normed space \(\mathrm{E}\) over \(\mathbf{R}\) ; let \(g\) be a continuous increasing real function on \(\mathrm{I}\) . Suppose that \(\mathbf{f}\) and \(g\) have right derivatives at all points of the relative complement in \([a, b]\) of a countable subset A of this interval (allowing \(g_r'(x)\) to be infinite at some of the points \(x \notin \mathbf{A}\) ), and suppose that at each of these points we have

\[
\left| \left| \mathbf {f} _ {d} ^ {\prime} (x) \right| \right| \leqslant g _ {r} ^ {\prime} (x).\tag{6}
\]

Under these hypotheses one has

\[
\| \mathbf {f} (b) - \mathbf {f} (a) \| \leqslant g (b) - g (a).\tag{7}
\]

The proof proceeds similarly to that of prop. 2. Let \(\varepsilon > 0\) be arbitrary, and \((a_{n})\) the sequence obtained by enumerating A in some order. Let J be the set of points \(y \in I\) such that, for all \(x\) such that \(a \leqslant x \leqslant y\) one has

\[
\| \mathbf {f} (x) - \mathbf {f} (a) \| \leqslant g (x) - g (a) + \varepsilon (x - a) + \varepsilon \sum_ {a _ {n} <   x} \frac {1}{2 ^ {n}};\tag{8}
\]

we shall show that J = I. One sees immediately, as in prop. 2, that J is an interval with left-hand endpoint a; if c is its right-hand endpoint then \(c \in J\) ; indeed, for all x \textless{} c one has (8), and a fortiori

\[
\| \mathbf {f} (x) - \mathbf {f} (a) \| \leqslant g (c) - g (a) + \varepsilon (c - a) + \varepsilon \sum_ {a _ {n} <   c} \frac {1}{2 ^ {n}}
\]

from which, letting \(x\) tend to \(c\) in this inequality, it follows from the continuity of \(\mathbf{f}\) that \(c\) satisfies (8).

Let us show that we must have c = b. So suppose that c \textless{} b and that moreover \(c \notin A\) : then \(\mathbf{f}_{d}^{\prime}(c)\) and \(g_{r}^{\prime}(c)\) exist and satisfy (6); suppose in the first place that \(g_{r}^{\prime}(c)\) (which is necessarily \(\geqslant 0\) ) is finite; then one can always write \(\mathbf{f}_{d}^{\prime}(c) = \mathbf{u}g_{r}^{\prime}(c)\) , with \(\|u\| \leqslant 1\) ; since the function \(\mathbf{f}(x) - \mathbf{u}g(x)\) has zero right derivative at the point c there must exist a y such that \(c < y \leqslant b\) and such that for \(c \leqslant x \leqslant y\) one has

\[
\| \mathbf {f} (x) - \mathbf {f} (c) - \mathbf {u} (g (x) - g (c)) \| \leqslant \varepsilon (x - c)
\]

from which

\[
\| \mathbf {f} (x) - \mathbf {f} (c) \| \leqslant g (x) - g (c) + \varepsilon (x - c)
\]

and, taking account of (8), in which x is replaced by c,

\[
\begin{array}{c} \| \mathbf {f} (x) - \mathbf {f} (a) \| \leqslant g (x) - g (a) + \varepsilon (x - a) + \varepsilon \sum_ {a _ {n} <   c} \frac {1}{2 ^ {n}} \\ \leqslant g (x) - g (a) + \varepsilon (x - a) + \varepsilon \sum_ {a _ {n} <   x} \frac {1}{2 ^ {n}}. \end{array}
\]

Thus one has \(y \in J\) , which is a contradiction. Suppose next that \(c \notin A\) and that \(g_{r}^{\prime}(c) = +\infty\) ; then there is a y such that \(c < y \leqslant b\) and such that for \(c \leqslant x \leqslant y\) one has on the one hand

\[
\left\| \mathbf {f} (x) - \mathbf {f} (c) \right\| \leqslant \left(\left\| \mathbf {f} _ {d} ^ {\prime} (c) \right\| + 1\right) (x - c)
\]

while on the other hand

\[
g (x) - g (c) \geqslant \left(\left\| \mathbf {f} _ {d} ^ {\prime} (c) \right\| + 1\right) (x - c)
\]

from which

\[
\| \mathbf {f} (x) - \mathbf {f} (c) \| \leqslant g (x) - g (c)
\]

and one concludes as above. Finally, if one has \(c = a_{k}\) , then there is a y such that \(c < y \leqslant b\) , and such that for \(c < x \leqslant y\) one has

\[
\| \mathbf {f} (x) - \mathbf {f} (c) \| \leqslant \frac {\varepsilon}{2 ^ {k}}
\]

from which, taking account of (8), with x replaced by c,

\[
\begin{array}{c} \| \mathbf {f} (x) - \mathbf {f} (a) \| \leqslant g (c) - g (a) + \varepsilon (c - a) + \varepsilon \sum_ {a _ {n} <   x} \frac {1}{2 ^ {n}} \\ \leqslant g (x) - g (a) + \varepsilon (x - a) + \varepsilon \sum_ {a _ {n} <   x} \frac {1}{2 ^ {n}} \end{array}
\]

which again entails a contradiction. The proof finishes as that of prop. 2.

Q.E.D.

Remarks. 1) Here again, in the statement of th. 2 one can replace the interval {[}a, b{[} by {]}a, b{]} and ``right derivative'' by ``left derivative''.

\begin{enumerate}
\def\labelenumi{\arabic{enumi})}
\setcounter{enumi}{1}
\tightlist
\item
  We shall show later how to identify the case of equality in (7), and also how to generalize th. 2 to the case where E is an arbitrary locally convex space, with the help of another method of proof which allows one to deduce th. 2 from th. 1.
\end{enumerate}

COROLLARY. For a continuous vector function on an interval \(I \subset R\) , with values in a normed space E over R, to be constant on I it suffices that it have zero right derivative at all points of the complement (relative to I) of a countable subset of I.

Remark. The proofs of ths. 1 and 2 rely in an essential manner on the special topological properties of the field R; one can give examples of valued fields K for which there are nonconstant linear maps of K to itself with zero derivative at every point (cf.~I, p.~37, exerc. 2).

PROPOSITION 3. Let f be a vector function with values in a normed space E over R, defined and continuous on an interval I ⊂ R, and right differentiable on the complement B (relative to I) of a countable subset of I; then for all points \(x_{0} \in B\) , \(x \in I\) , \(y \in I\) , one has (supposing that x \textless{} y, for example)

\[
\left\| \mathbf {f} (y) - \mathbf {f} (x) - \mathbf {f} _ {d} ^ {\prime} (x _ {0}) (y - x) \right\| \leqslant (y - x) \sup _ {z \in \mathrm{B}, x <   z <   y} \left\| \mathbf {f} _ {d} ^ {\prime} (z) - \mathbf {f} _ {d} ^ {\prime} (x _ {0}) \right\|.\tag{9}
\]

Indeed it suffices to apply th. 2, replacing f by the function

\[
\mathbf {f} (z) - \mathbf {f} _ {d} ^ {\prime} (x _ {0}) z,
\]

and \(g\) by the linear function whose derivative is \(\sup_{z\in \mathrm{B}, x < z < y}\left\| \mathbf{f}_d'(z) - \mathbf{f}_d'(x_0)\right\|\) .

Theorem 2 extends to vector functions of a complex variable:

PROPOSITION 4. Let \(\mathbf{f}\) be a continuous differentiable function of a complex variable defined on a convex open subset A of the field \(\mathbf{C}\) , with values in a normed space E over the field \(\mathbf{C}\) . If one has \(\| \mathbf{f}'(z)\| \leqslant m\) for all \(z\in \mathrm{A}\) , then one has \(\| \mathbf{f}(b) - \mathbf{f}(a)\| \leqslant m|b - a|\) for every pair of points \(a\) , \(b\) of A.

We put \(\mathbf{g}(t) = \frac{1}{b - a}\) \(\mathbf{f}(a + t(b - a))\) for \(0 \leqslant t \leqslant 1\) ; since \(\mathbf{g}'(t) = \mathbf{f}'(a + t(b - a))\) , applying th. 2 to the function \(\mathbf{g}\) yields the proposition immediately.

COROLLARY. For a vector function f of a complex variable, defined and continuous on an open set \(A \subset C\) , and with values in a normed space over C, to be constant, it suffices that it have zero derivative at every point of A.

Indeed, let a be an arbitrary point of A; the set B of points z of A where \(\mathbf{f}(z) = \mathbf{f}(a)\) is closed because f is continuous; it is also open, as is shown by applying prop. 4 (with m = 0) to a convex open neighbourhood, contained in A, of an arbitrary point of B; so is identical to A.

PROPOSITION 5. Let f be a vector function of a complex variable, defined, continuous and differentiable on a convex open set A ⊂ C, taking values in a normed space over the field C; then, no matter what the points \(x_{0}\) , x and y in A, one has

\[
\left| \left| \mathbf {f} (y) - \mathbf {f} (x) - \mathbf {f} _ {d} ^ {\prime} (x _ {0}) (y - x) \right| \right| \leqslant | y - x | \sup _ {z \in A} \left| \left| \mathbf {f} ^ {\prime} (z) - \mathbf {f} ^ {\prime} (x _ {0}) \right| \right|.\tag{10}
\]

It suffices to apply th. 2 to the function

\[
\mathbf {g} (t) = \mathbf {f} (x + t (y - x)) - \mathbf {f} ^ {\prime} (x _ {0}) (y - x) t
\]

on the interval {[}0, 1{]}.

